\begin{document}
\maketitle

\section{\textbf{Appointments}}
\cventry{Jan, 2019- present}{Postdoctoral scholar, Department of Bioengineering and Therapeutic Sciences}{University of California San Francisco}{San Francisco}{CA}{}

\section{\textbf{Education}}
\cventry{Sep, 2013- Dec, 2018}{Ph.D, Chemical Engineering}{University of California Santa Barbara}{Santa Barbara}{CA}{}

\cventry{Jan, 2014- Dec, 2016}{Graduate emphasis in Computational Science and Engineering}{University of California Santa Barbara}{Santa Barbara}{CA}{(Specialization with masters's equivalent coursework in numerical methods for ODEs and PDEs, and parallel computing)}

\cventry{Aug, 2008- May, 2013}{Masters and Bachelors of Technology (Hons.) integrated dual degree, Chemical Engineering}{Indian Institute of Technology Kharagpur}{Kharagpur}{WB, India}{}

\section{\textbf{Research Interests}}
\cvline{}{Integrative modeling of protein complexes, protein folding, mining relational protein structure data (contacts, crosslinks, coordination), probabilistic machine learning and Bayesian inference, molecular dynamics.}

\section{\textbf{Skills}}
\cvline{Softwares}{MATLAB, Mathematica, LAMMPS, GROMACS, AMBER, IMP, VMD, PyMol, UCSF Chimera}
\cvline{Frameworks}{Tensorflow-2,  PyMC3}
\cvline{Languages}{Python (\textasciitilde 5000 lines), Fortran-90 (\textasciitilde 1000 lines), C++ (\textasciitilde 1800 lines), Bash, TCL}
\cvline{Parallel-platforms}{MPI, OpenMP}

\section{\textbf{Research Experience}}
\cvline{\textbf{UCSF}}{Postdoc advisor: Andrej Sali}
\cvlistitem{\textbf{Integrative structural biology:} \newline Integrating fundamental (physics-based) constraints such as excluded volume and sequence connectivity requirements with mass-spectrometry derived residue-residue cross-links to refine experimental predictions of protein complexes implicated in the eukaryotic DNA replisome machinery.} 

\vspace{0.1in}
\cvlistitem{\textbf{Whole cell models:} \newline As a part of the Pancreatic Beta-Cell Consortium (collaborating across USC, UCLA, UCSF and SanghaiTech), working towards coupling traditional models of different parts, aspects, length and time-scales of beta cell biology, to build a "meta-model" of the whole-cell. Work involves developing computational toolchains for translating several different models and data like macroscopic pharmacokinetics, mesoscopic spatio-temporal Brownian dynamics simulations, microscopic molecular dynamics simulations of protein-protein interactions, biochemical pathways, soft-X-ray tomograms of insulin vesicles and other proteomic data into probabilistic graphical models as a rigorous underlying framework to couple such seemingly disparate models.}
\vspace{0.1in}

\cvline{\textbf{UCSB}}{PhD advisor: M. Scott Shell}
\cvlistitem{\textbf{Thesis: Next-Generation Coarse-Grained Models for Molecular Dynamics Simulations of Fluid Phase Equilibria and Protein Biophysics Using the Relative Entropy} \newline Introduced the local density potential as a manybody correction over traditional pair-potential based coarse grained forcefields for molecular dynamics simulations of liquid systems. Used local density augmented approaches to construct thermodynamically accurate coarse-grained simulations of liquid-liquid phase separation, examples of which are very very rare in the chemical engineering literature. Also developed orders of magnitude faster coarse grained models for polymer folding by removing water from such folding simulations but retaining the effect of water's hydrogen bonds using local density potentials.\newline \newline Developed hybrid physics + knowledge based algorithms for folding coarse-grained proteins (even very large proteins upto $\sim$ 247 residues) using molecular dynamics. Used simple coarse-grained protein models to investigate self assembly in (a) an engineering design problem that looks at the effects of confining surfaces and conjugated polymers on the secondary structures of surface tethered polypeptides (a putative model of micelles and polymer brushes) and (b) a fundamental (and open) structural biology problem: the connection between tertiary structure and amylodogenic propensity (leading to neurodegenerative diseases like Alzheimer's, Parkison's, ALS, etc.) in candidate peptide motifs.}
\vspace{0.1in}

\cvline{\textbf{IIT Kharagpur}}{Bachelors and Masters thesis advisor: Saikat Chakraborty}
\cvlistitem{\textbf{Master's Thesis: Formation and Stability of Mixing-Limited Patterns in Homogeneous Autocatalytic Reactors} \newline Developed a finite difference simulation of a coarse-grained plug-flow reactor to describe 2D transport limited patterns in candidate autocatalytic reactions. Developed bifurcation and linear perturbation theory approaches to determine parameters that give rise to meta-stable patterns.}

\vspace{0.1in}
\cvlistitem{\textbf{Bachelor's Thesis: Multiscale Analysis of Hypoxemia in Methemeglobin-Anemia.} \newline Developed a multi-component spatio-temporal reaction-diffusion model that couples the cellular (reaction), arterial (transport) and alveolar (mechanical) scales of human respiration, and applied it to characterize the role of nitric oxide in methemoglobin-anemia. Awarded best bachelors' thesis.}

\section{\textbf{Computational experience}}
\cvlistitem{\textbf{Bayesian inference in graphical models} - Experienced in rapidly prototyping inference problems for parameter selection in pharmacokinetic models, by building directed Bayesian networks describing the input-output correlations. Built customized distributions for PyMC3.}
\cvlistitem{\textbf{Analysis of high dimensional data} - Experienced in analyzing massively high-dimensional molecular dynamics trajectories to extract physically motivated low-dimensional projections and constructing meaningful visual representations.}
\cvlistitem{\textbf{Optimization} - Experienced in inverse design (of molecular mechanics interactions) using likelihood-maximization approaches}
\cvlistitem{\textbf{Parallel computing} - Parallelized runtime-critical parts of (PhD) group software using MPI}
\cvlistitem{\textbf{Object-oriented design} - Wrote coarse grained topology builders and analysis pipelines for protein structure modeling from scratch using a combination of Python and Fortran-90}
\cvlistitem{\textbf{Open-source contributions} - Wrote custom manybody potential energy function and post-processing pipelines for the molecular dynamics software LAMMPS.}
\cvlistitem{\textbf{Linux cluster management} - System administrator for (PhD) group Rocks 6.2 cluster. Installed Rocks and wrote post-install scripts to automate safe handling of the network-attached-storage and developed protocols to automate resource-sharing between lab members.}

\section{\textbf{Publications}}
%\cvlistitem{T. Sanyal*, J. Lu* and M.S. Shell, \textit{Polymer conjugation effects on secondary structure in tethered polypeptide assemblies}. Manuscript in preparation}
\vspace{0.1in}
\cvlistitem{T. Sanyal, J. Mittal and M.S. Shell,\textit{ A hybrid, bottom up, structurally-accurate, G$\bar{o}$-like coarse-grained protein model}. Journal of Chemical Physics (2019) 151, 044111 (featured article)}
\vspace{0.1in}
\cvlistitem{D. Rosenberger, T.Sanyal, M.S. Shell and N.F.A van der Vegt, \textit{Transferability of local density-assisted implicit solvation models for homogeneous fluid mixtures.} Journal of Chemical Theory and Computation (2019) 15, 2881-2895}
\vspace{0.1in}
\cvlistitem{M.P. Howard, W.F. Reinhart, T. Sanyal, M.S. Shell, A. Nikoubashman and A.Z. Panagiotopoulos, \textit{Evaporation-induced assembly of colloidal crystals.} Journal of Chemical Physics (2018) 149(\textbf{9})}
\vspace{0.1in}
\cvlistitem{T. Sanyal and M.S. Shell, \textit{Transferable coarse-grained models of liquid-liquid equilibrium using local density potentials optimized with the relative entropy.} Journal of Physical Chemistry-B (2018) 122 (\textbf{21}), 5678-5693}
\vspace{0.1in}
\cvlistitem{T. Sanyal and M.S. Shell, \textit{Coarse-grained models using local-density potentials optimized with the relative entropy: Application to implicit solvation.} Journal of Chemical Physics (2016) \textbf{145}, 034109}
\vspace{0.1in}
\cvlistitem{T. Sanyal and S. Chakraborty, \textit{Multiscale analysis of hypoxemia in methemoglobin-anemia.} Mathematical Biosciences (2013) \textbf{241}, 167-180}
\vspace{0.1in}
\cvlistitem{T. Sanyal and S. Chakraborty, \textit{Multiscale analysis of simultaneous uptake of two interacting gases in the human lungs and its application to methemoglobin anemia.} Computers and Chemical Engineering (2013) \textbf{59}, 226-242}
\vspace{0.1in}

\section{\textbf{Conference Presentations}}
\cvlistitem{Berkeley Mini Stat Mech Meeting; Berkeley, CA (Jan. 2019) \textbf{Poster}: \textit{ Coarse-grained models for protein folding and self-assembly with the relative entropy\\}}

\cvlistitem{10$^{\mathrm{th}}$ Annual Amgen Clorox Graduate Student Symposium; UC Santa Barbara, CA (Oct. 2017) \textbf{Talk}: \textit{A new approach to computationally fast coarse-grained models of hydrophobic interactions in fluids}\\}

\cvlistitem{AIChE Annual Meeting; San Francisco, CA (Nov. 2016) \textbf{Talk}: \textit{Improved coarse-grained models through local density potentials, optimized with the relative entropy}\\}

\cvlistitem{Berkeley Mini Stat Mech Meeting; Berkeley, CA (Jan. 2015) \textbf{Poster}: \textit{Building coarse grained models of solvation using local density dependent interactions with the relative entropy\\}}

\cvlistitem{Southern California Simulations in Science Conference; UC Santa Barbara, CA (Oct. 2015) \textbf{Poster}: \textit{Building coarse grained models of solvation using local density dependent interactions with the relative entropy}\\}

\cvlistitem{8$^{\mathrm{th}}$Amgen-Clorox Graduate Student Symposium; UC anta Barbara, CA (Sept. 2015) \textbf{Poster}: \textit{Robust models of coarse-grained interactions using multi-body potentials with the relative entropy}\\}

\cvlistitem{AIChE Annual Meeting; San Francisco, CA (Nov. 2013) \textbf{Poster}: \textit{Stability of mixing limited patterns in isothermal homogenous autocatalytic reactions}}

\section{\textbf{Teaching Assistance}}
\cvline{}{Undergraduate Thermodynamics (Fall 2018), Undergraduate Reaction Engineering (Spring 2017 ), Graduate Thermodynamics (Fall 2015), Introduction to Chemical Engineering(Fall 2014}

\section{\textbf{Undergraduates mentored}}
\cvlistitem{Jun Lu, UCSB, 2018-2019, Project: \textit{Effects of polymer conjugation on tethered peptides}}

\section{\textbf{Relevant Courses}}
\cvline{}{Thermodynamics, Equilibrium and Non-equilibrium statistical mechanics, Finite-Difference and Finite-Element methods for PDEs, Parallel computing}

\section{\textbf{Awards}}
\cvline{2012}{\textbf{Best bachelors thesis} IIT Kharagpur}
\cvline{2009}{\textbf{Technopreneur Promotion Programme, Ministry of Science and Technology, Govt. of India} - seed money for prototype scale-up of novel bio-reactors for algal bio-diesel generation}
\cvline{2009}{\textbf{Jagdis Bose National Science Talent Search (JBNSTS) foundation, WB, India} - fellowship}

\section{\textbf{Outreach Activities}}
\cvline{2014-2016}{\textbf{Co-founder and co-organizer, ChE Graduate Simulation Seminar Series}}
\cvlistitem{Conceptualized and started a seminar series with colleagues to address the unfulfilled need for a common collaborative and social platform for theory and computation researchers on campus.}
\cvlistitem{Developed extensive public outreach and attracted funding from ChE, Materials and Computer Science faculty, with participation from over 10 STEM fields and special attendances and keynotes from industry and national labs}

\cvline{2014-2015}{\textbf{President, India Association of Santa Barbara}}
\cvlistitem{Responsible for planning large-scale inclusive liberal arts and music festivals}
\cvlistitem{Increased participation from grad students and nearby Indian communities by 50 \%}
\end{document}